% Auto-generated complete chapter from Research/paquete_continuidad_agent_research.md
\chapter{paquete\_continuidad\_agent\_research}
\label{complete:research_065}

\begin{sourcemeta}
\textbf{Fuente:} \href{file:///E:/Laboral/Research/paquete_continuidad_agent_research.md}{Research/paquete\_continuidad\_agent\_research.md}\\
\textbf{Seccion:} research\\
\textbf{Estado:} research\_snapshot\\
\textbf{Rol:} Research / estudio de mercado\\
\textbf{Lineas originales:} 932\\
\textbf{Ultima modificacion:} 2026-06-11 16:50\\
\textbf{Deduplicacion conservadora:} 0 bloques repetidos omitidos; 0 caracteres repetidos no reimpresos.
\end{sourcemeta}

\section*{Resumen de orientacion}
Este documento permite continuar el proyecto en otro chat sin perder continuidad técnica.

\section*{Contenido completo depurado}
\addcontentsline{toc}{section}{Contenido completo depurado}




\subsection{0. Propósito del paquete}




Este documento permite continuar el proyecto en otro chat sin perder continuidad técnica.




El foco actual no es rehacer la investigación base, sino \textbf{continuar la fase de investigaciones con agente} para robustecer el roadmap internacional de Alexander Woodcock Salomón.




El flujo activo es:




\begin{mdcode}
Bloque de codigo: text
1. Ejecutar una tarea concreta con modo agente.
2. Pegar el resultado en el chat.
3. Auditar calidad, corregir vacíos y convertirlo en un módulo .md.
4. Pasar a la siguiente tarea.
5. Al final, integrar todo en una actualización del roadmap.
\end{mdcode}




\medskip\hrule\medskip




\subsection{1. Contexto base del proyecto}




\subsubsection{Perfil}




\begin{mdcode}
Bloque de codigo: text
Nombre: Alexander Woodcock Salomón
Ubicación: Colombia
Idioma nativo: Español
Inglés: C1 autopercibido
Alemán: inicial / no funcional
Portugués: estratégicamente relevante por ciudadanía portuguesa esperada
\end{mdcode}




\subsubsection{Posicionamiento profesional recomendado}




\begin{mdcode}
Bloque de codigo: text
Real-Time 3D Developer / Unity Technical Artist
focused on interactive technical visualization, CAD-to-realtime optimization, and Unity WebGL deployment.
\end{mdcode}




\subsubsection{Proyecto principal}




\begin{mdcode}
Bloque de codigo: text
TwinSight X500
Unity WebGL technical visualization prototype for inspecting and understanding the assembly of a Holybro X500 V2 drone.
\end{mdcode}




\subsubsection{Stack principal}




\begin{mdcode}
Bloque de codigo: text
Unity
C#
Unity WebGL
URP
UI Toolkit
Shader Graph
Blender
CAD-to-realtime optimization
Technical visualization
Interactive 3D
Python automation
Git/GitHub
\end{mdcode}




\subsubsection{Situación laboral}




\begin{mdcode}
Bloque de codigo: text
Formal industry experience in Unity / Technical Art: limited
Strongest proof: thesis / portfolio project
Target: remote international work from Colombia
Preferred model: contractor / B2B
Minimum initial target: USD 1.5k/month
Realistic 6-month target: USD 3k/month
Ambitious 12–24 month target: USD 6k/month
\end{mdcode}




\subsubsection{Movilidad}




\begin{mdcode}
Bloque de codigo: text
Current EU work authorization: No
Portuguese passport/citizenship: expected in approximately 2 years, considered highly likely by the user
Germany route: possible through legal marriage/residence route with German partner, but not current work authorization
Immediate strategy: remote from Colombia
Medium-term strategy: EU/Germany/Portugal once legally enabled
\end{mdcode}




\medskip\hrule\medskip




\subsection{2. Documentos ya generados}




El sistema principal ya llegó hasta:




\begin{mdcode}
Bloque de codigo: text
00–16: research, strategy, audit and planning
17–24: application materials and decision systems
25: spreadsheet tracker
26: weekly execution dashboard
27: conservative/balanced/aggressive scenarios
28: first agent benchmark module
\end{mdcode}




\subsubsection{Archivos importantes ya generados}




\begin{mdcode}
Bloque de codigo: text
17_cv_base_and_role_variants.md
18_linkedin_final_rewrite.md
19_github_readme_templates.md
20_portfolio_copy_and_site_structure.md
21_twinsight_demo_video_script_and_shotlist.md
22_outreach_templates_english_spanish.md
23_interview_answer_bank.md
24_offer_evaluation_scorecard.md
25_application_tracker_template.xlsx
26_weekly_execution_dashboard.md
27_strategy_scenarios_conservative_balanced_aggressive.md
28_competitor_profile_benchmark.md
\end{mdcode}




\subsubsection{Estado del archivo 28}




El archivo:




\begin{mdcode}
Bloque de codigo: text
28_competitor_profile_benchmark.md
\end{mdcode}




fue generado a partir de la tarea:




\begin{mdcode}
Bloque de codigo: text
A1_linkedin_profiles_benchmark
\end{mdcode}




Contiene:




\begin{itemize}
\item benchmark superficial de perfiles LinkedIn encontrados por snippets públicos;
\item familias de rol comparables;
\item headline patterns;
\item señales que Alexander debería copiar;
\item señales que no debería copiar;
\item implicaciones para CV, LinkedIn y portafolio.
\end{itemize}




\subsubsection{Limitación del archivo 28}




La búsqueda de LinkedIn fue limitada porque muchas páginas requieren login. Por tanto, el archivo 28 debe considerarse:




\begin{mdcode}
Bloque de codigo: text
surface-level LinkedIn positioning benchmark
\end{mdcode}




No debe considerarse un análisis completo de perfiles.




\medskip\hrule\medskip




\subsection{3. Estado actual de tareas con agente}




\subsubsection{Completado}




\begin{center}
\scriptsize
\begin{longtable}{>{\raggedright\arraybackslash}p{0.225\linewidth} >{\raggedright\arraybackslash}p{0.225\linewidth} >{\raggedright\arraybackslash}p{0.225\linewidth} >{\raggedright\arraybackslash}p{0.225\linewidth}}
\toprule
Código & Tarea & Estado & Salida \\
\midrule
\endfirsthead
\toprule
Código & Tarea & Estado & Salida \\
\midrule
\endhead
A1 & LinkedIn profiles benchmark & Completada & \texttt{28\_competitor\_profile\_benchmark.md} \\
\bottomrule
\end{longtable}
\end{center}




\subsubsection{Intentada pero incompleta}




\begin{center}
\scriptsize
\begin{longtable}{>{\raggedright\arraybackslash}p{0.225\linewidth} >{\raggedright\arraybackslash}p{0.225\linewidth} >{\raggedright\arraybackslash}p{0.225\linewidth} >{\raggedright\arraybackslash}p{0.225\linewidth}}
\toprule
Código & Tarea & Estado & Problema \\
\midrule
\endfirsthead
\toprule
Código & Tarea & Estado & Problema \\
\midrule
\endhead
A2 & Public CV / resume benchmark & Incompleta & El agente encontró pocas fuentes y terminó diciendo que no podía recopilar más \\
\bottomrule
\end{longtable}
\end{center}




\subsubsection{Próxima acción recomendada}




Reintentar A2 con un prompt más robusto y menos restrictivo.




Razón:




\begin{itemize}
\item la primera ejecución de A2 fue insuficiente;
\item sí encontró fuentes útiles, pero no las explotó bien;
\item se deben permitir fuentes como portfolios personales, PDFs, resume databases, university samples, ArtStation resume pages y páginas personales;
\item se debe pedir mínimo 12–20 ejemplos, no 30–50;
\item debe aceptar ejemplos anónimos o sample resumes siempre que sean concretos.
\end{itemize}




\medskip\hrule\medskip




\subsection{4. Regla de uso de agente}




El agente se usa para:




\begin{mdcode}
Bloque de codigo: text
public web collection
structured source data
live verification
tables
URLs
date checked
\end{mdcode}




No se usa para:




\begin{mdcode}
Bloque de codigo: text
final strategy
CV writing
LinkedIn writing
portfolio copy
negotiation judgment
legal interpretation
narrative synthesis
\end{mdcode}




La síntesis final la hace el asistente después de recibir los datos.




\medskip\hrule\medskip




\subsection{5. Secuencia general de nuevos módulos}




Estos son los módulos adicionales propuestos después del sistema 00–27.




\begin{mdcode}
Bloque de codigo: text
28_competitor_profile_benchmark.md
29_portfolio_case_study_benchmark.md
30_live_job_postings_market_snapshot.md
31_verified_company_targets_expanded.md
32_education_scholarships_masters_fellowships.md
33_competitions_awards_hackathons_thesis_opportunities.md
34_volunteering_open_source_experience_pipeline.md
35_recruiter_outreach_database.md
36_content_strategy_linkedin_artstation_github.md
37_salary_contracting_legal_update.md
38_agent_research_synthesis_and_roadmap_update.md
\end{mdcode}




\subsubsection{Ajuste necesario}




Como \texttt{28\_competitor\_profile\_benchmark.md} ya fue generado solamente con A1, las próximas tareas de benchmark profesional pueden manejarse de dos formas:




\paragraph{Opción recomendada}



Crear archivos complementarios:




\begin{mdcode}
Bloque de codigo: text
28B_public_cv_resume_benchmark.md
28C_github_profiles_benchmark.md
28D_artstation_portfolio_benchmark.md
\end{mdcode}




Luego crear:




\begin{mdcode}
Bloque de codigo: text
28E_competitor_benchmark_synthesis.md
\end{mdcode}




\paragraph{Opción alternativa}



Regenerar el archivo 28 como versión completa:




\begin{mdcode}
Bloque de codigo: text
28_competitor_profile_benchmark_v2.md
\end{mdcode}




La opción recomendada es \textbf{28B / 28C / 28D / 28E}, porque evita sobreescribir trabajo ya hecho.




\medskip\hrule\medskip




\subsection{6. Orden inmediato recomendado}




\subsubsection{Próximas tareas}




\begin{center}
\scriptsize
\begin{longtable}{>{\raggedright\arraybackslash}p{0.184\linewidth} >{\raggedright\arraybackslash}p{0.184\linewidth} >{\raggedright\arraybackslash}p{0.184\linewidth} >{\raggedright\arraybackslash}p{0.184\linewidth} >{\raggedright\arraybackslash}p{0.184\linewidth}}
\toprule
Orden & Código & Tarea & Usar agente & Deep Research \\
\midrule
\endfirsthead
\toprule
Orden & Código & Tarea & Usar agente & Deep Research \\
\midrule
\endhead
1 & A2-Retry & Public CV / resume benchmark & Sí & No \\
2 & A3 & GitHub profiles benchmark & Sí & No \\
3 & A4 & ArtStation / portfolio benchmark & Sí & No \\
4 & A5 & Technical visualization case studies & Sí & No \\
5 & F1 & Demo reels benchmark & Sí & No \\
6 & F3 & ArtStation breakdown benchmark & Sí & No \\
7 & B1 & Unity / Real-Time 3D live jobs & Sí & No \\
8 & B2 & Digital Twin / Simulation live jobs & Sí & No \\
9 & B5 & LATAM nearshore jobs & Sí & No \\
10 & B6 & Top 30 company verification & Sí & No \\
\bottomrule
\end{longtable}
\end{center}




\subsubsection{Motivo del orden}




Primero se mejora el benchmark de presentación profesional:




\begin{mdcode}
Bloque de codigo: text
CV → GitHub → ArtStation/Portfolio → Case Studies → Demo reels
\end{mdcode}




Luego se pasa al mercado vivo:




\begin{mdcode}
Bloque de codigo: text
jobs → companies → recruiters → education → competitions → volunteering
\end{mdcode}




Esto evita gastar usos de agente en ofertas que pueden caducar antes de que los materiales estén listos.




\medskip\hrule\medskip




\subsection{7. Prompt recomendado para reintentar A2}




Usar este prompt en el próximo chat o en el mismo chat, con modo agente activo.




\begin{mdcode}
Bloque de codigo: text
Execute A2_retry_public_cv_resume_benchmark.

Context:
We already attempted A2, but the result was incomplete. This retry should be more flexible and collect concrete public resume/CV examples from personal websites, PDF resumes, resume databases, university sample packets, GitHub resume pages, ArtStation resume pages, portfolio CV pages, and public blind CVs.

Objective:
Collect 12–20 public CV/resume examples relevant to Alexander’s target positioning:
Real-Time 3D Developer / Unity Technical Artist / Unity WebGL Developer / Technical Visualization Developer / XR Developer / Simulation Developer / Tools-Pipeline Developer.

Acceptable source types:
- real candidate public PDF resumes;
- personal website resume pages;
- portfolio pages with resume section;
- GitHub resume READMEs;
- ArtStation resume pages;
- blind CVs from reputable recruiting databases;
- university sample resumes if they contain concrete examples relevant to game development, technical art, creative technology, computational design, Unity, XR, simulation, visualization, or technical tools.

Do not use generic advice articles unless they include a concrete resume sample.

Output structured tables only.

Fields:
- Source label
- Resume/CV URL
- Source type
- Role target
- Country/region if visible
- Seniority
- Resume structure
- Headline/title used
- Summary style
- Skills organization
- Project/experience structure
- Metrics used
- Education wording
- Portfolio/GitHub links if visible
- What Alexander can adapt
- What Alexander should not copy
- Confidence
- Date checked

Rules:
- Use only publicly available information.
- Do not invent missing data.
- If unavailable, write “not found”.
- Keep cells short.
- Include source URLs.
- Tables only.
- Do not write strategy.
- Do not write narrative recommendations.
- Mark generated/sample resumes as “sample”, not real candidate.
\end{mdcode}




\subsubsection{Salida esperada después de A2-Retry}




Cuando el agente entregue resultados, generar:




\begin{mdcode}
Bloque de codigo: text
28B_public_cv_resume_benchmark.md
\end{mdcode}




Ese archivo debe contener:




\begin{itemize}
\item fuentes recolectadas;
\item patrones de estructura de CV;
\item qué adaptar al CV de Alexander;
\item qué evitar;
\item actualización concreta para \texttt{17\_cv\_base\_and\_role\_variants.md};
\item prompt para la siguiente tarea A3.
\end{itemize}




\medskip\hrule\medskip




\subsection{8. Prompt para A3 — GitHub profiles benchmark}




Usar después de cerrar A2-Retry.




\begin{mdcode}
Bloque de codigo: text
Execute A3_github_profiles_benchmark.

Objective:
Collect public GitHub profile examples relevant to Alexander’s positioning:
Real-Time 3D Developer, Unity Technical Artist, Unity WebGL Developer, Technical Visualization Developer, XR Developer, Simulation Developer, Tools/Pipeline Developer, Python automation/tooling developer.

Find 20–35 GitHub profiles or repositories.

Prioritize:
- Unity projects with strong README;
- WebGL / Three.js / WebGPU projects;
- technical visualization projects;
- simulation or digital twin projects;
- Unity tools/editor tools;
- technical artist tools;
- Python automation / LLM tooling repos;
- portfolios using GitHub Pages.

Output structured tables only.

Fields:
- GitHub profile/repo label
- URL
- Source type: profile / repo / GitHub Pages / organization
- Role relevance
- Main technologies
- README structure
- Pinned repos if profile
- Project media used
- Badges/topics used
- Demo links
- Portfolio links
- License/disclaimer handling
- Documentation quality
- What Alexander can adapt
- What Alexander should not copy
- Confidence
- Date checked

Rules:
- Use only public GitHub data.
- Do not clone private repos.
- Do not invent unavailable data.
- If unavailable, write “not found”.
- Tables only.
- No narrative strategy.
- Include source URLs.
\end{mdcode}




Expected output file:




\begin{mdcode}
Bloque de codigo: text
28C_github_profiles_benchmark.md
\end{mdcode}




\medskip\hrule\medskip




\subsection{9. Prompt para A4 — ArtStation / portfolio benchmark}




\begin{mdcode}
Bloque de codigo: text
Execute A4_artstation_portfolio_benchmark.

Objective:
Collect public ArtStation and portfolio examples relevant to Alexander’s technical art and real-time 3D positioning.

Find 20–35 examples across:
- Unity Technical Artist portfolios;
- Technical Artist portfolios;
- shader breakdowns;
- VFX breakdowns;
- real-time 3D portfolios;
- 3D optimization breakdowns;
- Blender technical breakdowns;
- character/portrait technical breakdowns;
- XR / interactive 3D portfolios;
- technical visualization portfolios.

Output structured tables only.

Fields:
- Artist / profile label
- URL
- Platform
- Role target
- Seniority if visible
- Project type
- Breakdown structure
- Images used
- Videos used
- Technical text depth
- Tags/skills
- Tools shown
- Portfolio/GitHub links if visible
- What Alexander can adapt
- What Alexander should not copy
- Confidence
- Date checked

Rules:
- Use only public pages.
- Do not invent missing data.
- If unavailable, write “not found”.
- Tables only.
- No narrative strategy.
- Include source URLs.
\end{mdcode}




Expected output file:




\begin{mdcode}
Bloque de codigo: text
28D_artstation_portfolio_benchmark.md
\end{mdcode}




\medskip\hrule\medskip




\subsection{10. Prompt para A5 — Technical visualization case studies}




\begin{mdcode}
Bloque de codigo: text
Execute A5_technical_visualization_case_studies.

Objective:
Collect public case studies relevant to TwinSight X500 and Alexander’s positioning:
Unity WebGL technical visualization, CAD-to-realtime, digital twin visualization, product configurators, simulation visualization, XR training, industrial visualization, technical inspection tools.

Find 20–30 examples from:
- company case studies;
- portfolio case studies;
- Unity showcases;
- WebGL/Three.js showcases;
- digital twin demos;
- industrial/product visualization studios;
- XR simulation demos.

Output structured tables only.

Fields:
- Case study title
- URL
- Company/person
- Industry
- Technology stack if visible
- Project type
- Problem statement
- Solution summary
- Features shown
- Media used
- Metrics/results if visible
- Case study structure
- Relevance to TwinSight
- What Alexander can adapt
- What Alexander should not copy
- Confidence
- Date checked

Rules:
- Use public sources only.
- Prefer official case studies or direct portfolio pages.
- Do not invent stack or metrics.
- If unavailable, write “not found”.
- Tables only.
- No strategy.
- Include source URLs.
\end{mdcode}




Expected output file:




\begin{mdcode}
Bloque de codigo: text
29_portfolio_case_study_benchmark.md
\end{mdcode}




\medskip\hrule\medskip




\subsection{11. Prompt para F1 — Demo reels benchmark}




\begin{mdcode}
Bloque de codigo: text
Execute F1_demo_reels_benchmark.

Objective:
Collect public demo reel examples relevant to TwinSight X500:
technical artist demo reels, Unity technical reels, real-time 3D developer reels, technical visualization demos, XR/simulation demos, CAD-to-realtime/product visualization demos.

Find 15–25 examples.

Output structured tables only.

Fields:
- Demo label
- URL
- Platform
- Role target
- Duration
- Opening structure
- Shot sequence
- Captions/text usage
- Voiceover yes/no
- Technical breakdown yes/no
- Tools shown
- Metrics shown
- Call-to-action
- What Alexander can adapt
- What Alexander should not copy
- Confidence
- Date checked

Rules:
- Use only public videos/pages.
- Do not invent duration if not visible.
- If unavailable, write “not found”.
- Tables only.
- Include URLs.
\end{mdcode}




Expected output file:




\begin{mdcode}
Bloque de codigo: text
29B_demo_reels_benchmark.md
\end{mdcode}




\medskip\hrule\medskip




\subsection{12. Prompt para F3 — ArtStation breakdown benchmark}




\begin{mdcode}
Bloque de codigo: text
Execute F3_artstation_breakdown_benchmark.

Objective:
Collect ArtStation or portfolio breakdowns useful for Alexander’s Blender portrait and technical art presentation.

Find 15–25 examples involving:
- topology breakdown;
- material breakdown;
- grooming breakdown;
- lighting breakdown;
- shader breakdown;
- real-time asset breakdown;
- before/after optimization;
- character study presentation;
- technical art process documentation.

Output structured tables only.

Fields:
- Project label
- URL
- Platform
- Breakdown type
- Tools shown
- Images used
- Video/GIF used
- Technical text depth
- Tags
- What Alexander can adapt
- What Alexander should not copy
- Confidence
- Date checked

Rules:
- Public pages only.
- No invented data.
- If unavailable, write “not found”.
- Tables only.
- Include URLs.
\end{mdcode}




Expected output file:




\begin{mdcode}
Bloque de codigo: text
29C_artstation_breakdown_benchmark.md
\end{mdcode}




\medskip\hrule\medskip




\subsection{13. Mercado vivo: tareas posteriores}




\subsubsection{B1 — Unity / Real-Time 3D live jobs}




Expected output:




\begin{mdcode}
Bloque de codigo: text
30_live_job_postings_market_snapshot.md
\end{mdcode}




Prompt resumido:




\begin{mdcode}
Bloque de codigo: text
Collect live job postings for Unity Developer, Real-Time 3D Developer, Unity Technical Artist, Unity WebGL Developer, Interactive 3D Developer.

Fields:
company, role, URL, date checked, region, remote scope, salary if visible, seniority, requirements, portfolio expectations, work authorization, contractor possibility, relevance to Alexander, confidence.

Tables only. No strategy. Do not assume remote means worldwide.
\end{mdcode}




\subsubsection{B2 — Digital Twin / Simulation live jobs}




Output can append to:




\begin{mdcode}
Bloque de codigo: text
30_live_job_postings_market_snapshot.md
\end{mdcode}




\subsubsection{B5 — LATAM nearshore live jobs}




Output can append to:




\begin{mdcode}
Bloque de codigo: text
30_live_job_postings_market_snapshot.md
\end{mdcode}




\subsubsection{B6 — Top 30 company verification}




Expected output:




\begin{mdcode}
Bloque de codigo: text
31_verified_company_targets_expanded.md
\end{mdcode}




\medskip\hrule\medskip




\subsection{14. Educación y oportunidades}




Future tasks:




\begin{mdcode}
Bloque de codigo: text
C1_high_roi_short_courses
C2_rebelway_deep_audit
C3_masters_computer_graphics_xr_simulation
C4_scholarships_portugal_germany_eu
C5_phd_research_assistantships
D1_thesis_innovation_competitions
D2_conferences_events_demo_venues
D3_paper_poster_publication_venues
E1_technical_volunteering_opportunities
E2_open_source_unity_webgl_projects
E3_remote_internships_apprenticeships
E4_aligned_freelance_gig_opportunities
\end{mdcode}




Expected modules:




\begin{mdcode}
Bloque de codigo: text
32_education_scholarships_masters_fellowships.md
33_competitions_awards_hackathons_thesis_opportunities.md
34_volunteering_open_source_experience_pipeline.md
\end{mdcode}




\medskip\hrule\medskip




\subsection{15. Síntesis final después de todas las tareas}




Cuando estén listos los módulos 28B–37, generar:




\begin{mdcode}
Bloque de codigo: text
38_agent_research_synthesis_and_roadmap_update.md
\end{mdcode}




Ese módulo debe responder:




\begin{mdcode}
Bloque de codigo: text
1. Qué cambió frente al roadmap original.
2. Qué no cambió.
3. Qué empresas suben a prioridad A.
4. Qué roles bajan de prioridad.
5. Qué cambios aplicar a CV.
6. Qué cambios aplicar a LinkedIn.
7. Qué cambios aplicar a GitHub.
8. Qué cambios aplicar al portafolio.
9. Qué cursos o programas sí justifican tiempo.
10. Qué becas o maestrías entran al roadmap 12–24 meses.
11. Qué concursos o venues valen la pena.
12. Qué oportunidades de experiencia práctica son viables.
13. Qué ajustes hacer al tracker.
14. Qué cambia en el plan 30/60/90.
\end{mdcode}




\medskip\hrule\medskip




\subsection{16. Criterios de calidad para resultados de agente}




Un resultado de agente es útil si tiene:




\begin{mdcode}
Bloque de codigo: text
URLs verificables
fecha de revisión
fuente concreta
datos estructurados
campos “not found” cuando no hay información
distinción entre perfil real y sample
confianza por fila
celdas cortas
sin narrativa innecesaria
\end{mdcode}




Un resultado de agente es insuficiente si:




\begin{mdcode}
Bloque de codigo: text
termina diciendo que no puede buscar
trae solo 2–3 ejemplos
trae artículos genéricos sin ejemplos
mezcla estrategia con datos
inventa información no visible
no incluye URLs
no incluye fecha
no separa muestras reales de samples
\end{mdcode}




Cuando un resultado es insuficiente:




\begin{mdcode}
Bloque de codigo: text
1. No generar módulo final todavía, o generar módulo marcándolo como incompleto.
2. Reintentar con prompt más específico.
3. Reducir volumen pedido.
4. Ampliar fuentes aceptables.
5. Pedir salida parcial.
\end{mdcode}




\medskip\hrule\medskip




\subsection{17. Estado exacto para continuar}




\subsubsection{Última acción completada correctamente}




\begin{mdcode}
Bloque de codigo: text
A1_linkedin_profiles_benchmark completada.
28_competitor_profile_benchmark.md generado.
\end{mdcode}




\subsubsection{Última acción fallida o incompleta}




\begin{mdcode}
Bloque de codigo: text
A2_public_cv_resume_benchmark intentada, pero insuficiente.
\end{mdcode}




\subsubsection{Próxima acción exacta}




\begin{mdcode}
Bloque de codigo: text
Ejecutar A2_retry_public_cv_resume_benchmark con modo agente activado.
\end{mdcode}




\subsubsection{Después de recibir resultados A2-Retry}




Generar:




\begin{mdcode}
Bloque de codigo: text
28B_public_cv_resume_benchmark.md
\end{mdcode}




Luego indicar:




\begin{mdcode}
Bloque de codigo: text
Siguiente tarea: A3_github_profiles_benchmark
Modo agente: Sí
Deep Research: No
\end{mdcode}




\medskip\hrule\medskip




\subsection{18. Instrucción para el próximo chat}




Copiar este bloque al iniciar un chat nuevo:




\begin{mdcode}
Bloque de codigo: text
Continuamos el proyecto Laboral de Alexander Woodcock Salomón.

No rehacer la estrategia base. Ya existen módulos 00–27, el tracker 25 y el módulo 28 de benchmark LinkedIn.

Estado actual:
- A1_linkedin_profiles_benchmark completado.
- 28_competitor_profile_benchmark.md generado.
- A2_public_cv_resume_benchmark fue intentado pero quedó incompleto.
- Próxima tarea: ejecutar A2_retry_public_cv_resume_benchmark con agente.
- Después de recibir los datos, generar 28B_public_cv_resume_benchmark.md.
- Mantener el flujo: agente recolecta datos → asistente audita/sintetiza → genera .md → propone siguiente tarea.

Usar la estrategia central:
Real-Time 3D Developer / Unity Technical Artist focused on Unity WebGL, CAD-to-realtime optimization and technical visualization.

No usar Deep Research salvo que se pida explícitamente para investigación profunda. Para las próximas tareas usar modo agente, no Deep Research.
\end{mdcode}




\medskip\hrule\medskip




\subsection{19. Prompt compacto para continuar de inmediato}




Usar este prompt como siguiente paso:




\begin{mdcode}
Bloque de codigo: text
Execute A2_retry_public_cv_resume_benchmark.

Context:
A2 was attempted but incomplete. Retry with broader source types and lower target volume.

Objective:
Collect 12–20 concrete public resume/CV examples relevant to Real-Time 3D Developer, Unity Technical Artist, Unity WebGL Developer, Technical Visualization Developer, XR Developer, Simulation Developer, Tools/Pipeline Developer.

Accept:
- real public PDF resumes;
- personal website resume pages;
- portfolio CV pages;
- GitHub resume READMEs;
- ArtStation resume pages;
- blind CVs from recruiting databases;
- university sample resumes if they include concrete examples relevant to Unity, game development, technical art, creative technology, computational design, XR, simulation, visualization, or technical tools.

Output structured tables only.

Fields:
- Source label
- Resume/CV URL
- Source type
- Role target
- Country/region if visible
- Seniority
- Resume structure
- Headline/title used
- Summary style
- Skills organization
- Project/experience structure
- Metrics used
- Education wording
- Portfolio/GitHub links if visible
- What Alexander can adapt
- What Alexander should not copy
- Confidence
- Date checked

Rules:
- Public sources only.
- Do not invent data.
- If unavailable, write “not found”.
- Keep cells short.
- Include source URLs.
- Tables only.
- No strategy.
- Mark generated/sample resumes as sample, not real candidate.
\end{mdcode}




\medskip\hrule\medskip




\subsection{20. Prioridad real}




El proyecto ya tiene suficiente estrategia.




La prioridad ahora es:




\begin{mdcode}
Bloque de codigo: text
1. mejorar evidencia pública;
2. comparar contra perfiles reales;
3. aplicar cambios a CV / LinkedIn / GitHub / portfolio;
4. empezar aplicaciones con tracker;
5. usar agente para verificar oportunidades vivas.
\end{mdcode}




No volver a abrir una investigación general sin razón.



